% Compile with pdfLaTeX: latexmk -pdf Tianjian_Qin_CV.tex
% The document uses fonts and packages supplied by TeX Live.
\documentclass[10pt,a4paper]{article}
\usepackage[left=1.7cm,right=1.7cm,top=1.45cm,bottom=1.45cm]{geometry}
\usepackage[T1]{fontenc}
\usepackage[utf8]{inputenc}
\usepackage[default]{lato}
\usepackage{microtype}
\usepackage{anyfontsize}
\usepackage{xcolor}
\usepackage{tabularx,array}
\newcolumntype{Y}{>{\RaggedRight\arraybackslash}X}
\usepackage{enumitem}
\usepackage{ragged2e}
\usepackage{needspace}
\usepackage{fancyhdr}
\usepackage{lastpage}
\usepackage[most]{tcolorbox}
\usepackage{xurl}
\usepackage{fontawesome5}
\usepackage{hyperref}
\input{glyphtounicode}
\pdfgentounicode=1

\definecolor{Ink}{HTML}{17263B}
\definecolor{Accent}{HTML}{315BA6}
\definecolor{Muted}{HTML}{566273}
\definecolor{Rule}{HTML}{CFD8E5}
\definecolor{Tint}{HTML}{F0F4FA}
\definecolor{Warm}{HTML}{8D6539}

\hypersetup{
  colorlinks=true,urlcolor=Accent,linkcolor=Accent,
  pdfauthor={Tianjian Qin},
  pdftitle={Tianjian Qin | Computational Biology and AI/ML},
  pdfsubject={Academic curriculum vitae},
  pdfkeywords={Computational biology, AI, machine learning, infectious disease modelling,
    network science, phylodynamics, One Health, research software},
  pdfdisplaydoctitle=true,
  pdflang={en-GB}
}
\urlstyle{same}
\setlength{\parindent}{0pt}
\setlength{\parskip}{0pt}
\setlength{\tabcolsep}{0pt}
\setlength{\footskip}{8mm}
\setlength{\emergencystretch}{2em}
\hyphenpenalty=10000
\exhyphenpenalty=10000
\raggedbottom

\pagestyle{fancy}
\fancyhf{}
\renewcommand{\headrulewidth}{0pt}
\renewcommand{\footrulewidth}{0.4pt}
\renewcommand{\footrule}{\hbox to\headwidth{\color{Rule}\leaders\hrule height\footrulewidth\hfill}\vskip 4pt}
\fancyfoot[L]{\fontsize{8}{10}\selectfont\color{Muted}Tianjian Qin\enspace\textcolor{Rule}{|}\enspace Curriculum vitae}
\fancyfoot[C]{\fontsize{8}{10}\selectfont\color{Muted}September 2026}
\fancyfoot[R]{\fontsize{8}{10}\selectfont\color{Ink}\thepage\enspace/\enspace\pageref*{LastPage}}

\newcommand{\sep}{\enspace\textcolor{Rule}{|}\enspace}
\newcommand{\middotsep}{\enspace\textcolor{Muted}{\textperiodcentered}\enspace}
\newcommand{\smallcapslabel}[1]{{\fontsize{8.3}{10}\selectfont\bfseries\color{Accent}\textls[105]{\MakeUppercase{#1}}}}
\newcounter{cvsectioncount}
\newcommand{\sectionhead}[1]{%
  \Needspace{4\baselineskip}\par\vspace{9pt}
  \stepcounter{cvsectioncount}\pdfbookmark[1]{#1}{section.\thecvsectioncount}
  {\fontsize{10.3}{12}\selectfont\bfseries\color{Ink}\textls[95]{\MakeUppercase{#1}}}\par
  \vspace{4pt}{\color{Rule}\hrule height 0.5pt}\vspace{7pt}
}
\newcommand{\compactsectionhead}[1]{%
  \Needspace{2.2\baselineskip}\par\vspace{5pt}
  {\fontsize{9.9}{11.5}\selectfont\bfseries\color{Ink}\textls[90]{\MakeUppercase{#1}}}\par
  \vspace{2.5pt}{\color{Rule}\hrule height 0.5pt}\vspace{4.5pt}
}
\newcommand{\subhead}[1]{%
  \Needspace{2.2\baselineskip}\par\vspace{7pt}
  {\fontsize{9.2}{11}\selectfont\bfseries\color{Accent}#1}\par\vspace{4pt}
}
\newcommand{\pageheading}[2]{%
  \pdfbookmark[0]{#2}{page.\thepage}%
  \smallcapslabel{Tianjian Qin}\hfill{\fontsize{8.3}{10}\selectfont\color{Muted}#1}\par
  \vspace{7pt}{\fontsize{23}{26}\selectfont\bfseries\color{Ink}#2}\par
  \vspace{9pt}{\color{Accent}\hrule height 1.5pt}\vspace{3pt}
}
\newcommand{\datedentry}[4]{%
  \Needspace{2.2\baselineskip}\par
  \begin{tabularx}{\linewidth}{@{}>{\RaggedRight\arraybackslash}p{2.65cm}@{\hspace{0.35cm}}Y@{}}
    {\fontsize{9.2}{12}\selectfont\bfseries\color{Accent}#1} &
    {\bfseries\color{Ink}#2}\par
    {\fontsize{9.4}{12}\selectfont\color{Muted}#3}\par
    \vspace{2pt}#4
  \end{tabularx}\par\vspace{7pt}
}
\newcommand{\proof}[4]{%
  \Needspace{5\baselineskip}\par
  {\bfseries\color{Ink}#1}\hfill{\fontsize{8.7}{11}\selectfont\color{Accent}#2}\par
  \vspace{2pt}#3\par
  \vspace{3pt}{\fontsize{8.9}{11}\selectfont\color{Muted}#4}\par\vspace{8pt}
}
\newcommand{\programmetag}[2]{%
  {\fontsize{10.2}{12.4}\selectfont\bfseries\color{Accent}\MakeUppercase{#1}%
  \hspace{0.4em}\textcolor{Rule}{|}\hspace{0.4em}#2}%
}
\newcommand{\programme}[5]{%
  \Needspace{3.2\baselineskip}\par
  \begin{tabularx}{\linewidth}{@{}Y@{\hspace{0.25cm}}>{\RaggedLeft\arraybackslash}p{4.45cm}@{}}
    {\bfseries\color{Ink}#1} & \programmetag{#2}{#3}
  \end{tabularx}\par
  {\fontsize{8.8}{10.5}\selectfont\color{Muted}#4}\par
  \vspace{0.7pt}{\fontsize{9.45}{11.3}\selectfont #5}\par\vspace{2pt}
}
\newcommand{\skillrow}[2]{\textbf{\color{Ink}#1} & #2\\[2.5pt]}
\newcommand{\qname}{\textbf{Qin, T.}}
\newcommand{\qjname}{\textbf{Qin, T.-J.}}
\newcommand{\senior}{\textsuperscript{\dag}}
\newcommand{\doilink}[1]{\href{https://doi.org/#1}{\textcolor{Accent}{DOI}}}
\newcommand{\contacticon}[1]{\textcolor{Accent}{\faIcon{#1}}\hspace{0.32em}}
\newcommand{\expertiseicon}[1]{\makebox[1.35em][c]{\textcolor{Accent}{\faIcon{#1}}}}
\newcommand{\skillrowicon}[3]{\expertiseicon{#1}\textbf{\color{Ink}#2} & #3\\[2.5pt]}

\newlist{cvitems}{itemize}{1}
\setlist[cvitems]{label=\textcolor{Accent}{\textbullet},leftmargin=11pt,
  labelsep=5pt,topsep=2pt,itemsep=2pt,parsep=0pt,partopsep=0pt}
\newlist{pubs}{enumerate}{1}
\setlist[pubs]{label=\textcolor{Accent}{\bfseries\arabic*.},leftmargin=17pt,
  labelsep=6pt,topsep=1pt,itemsep=4.8pt,parsep=0pt,partopsep=0pt}

\newtcolorbox{showcase}{
  enhanced,colback=Tint,colframe=Tint,boxrule=0pt,sharp corners,
  borderline west={2pt}{0pt}{Accent},
  left=10pt,right=10pt,top=8pt,bottom=8pt,
  before skip=0pt,after skip=3pt
}

\begin{document}
\fontsize{10.2}{13.1}\selectfont
\color{Ink}\RaggedRight

% Professional identity and selected evidence.
\pdfbookmark[0]{Profile and research experience}{overview}
{\fontsize{32}{35}\selectfont\bfseries Tianjian Qin}{\fontsize{17}{20}\selectfont\color{Muted}\enspace PhD}\hfill
{\fontsize{9}{11}\selectfont\color{Muted}\contacticon{map-marker-alt}Netherlands}\par
\vspace{5pt}
{\fontsize{15}{18}\selectfont\color{Accent}Computational Biologist \& AI/ML Researcher}\par
\vspace{4pt}
{\fontsize{10}{13}\selectfont Postdoctoral Researcher\middotsep Wageningen University \& Research}\par
\vspace{9pt}
{\fontsize{9.1}{12}\selectfont
  \contacticon{envelope}\href{mailto:tianjian.qin@wur.nl}{tianjian.qin@wur.nl}\sep
  \contacticon{globe}\href{https://qtj.me}{qtj.me}\sep
  \contacticon{project-diagram}\href{https://herdlink.nl}{herdlink.nl}\sep
  \contacticon{github}\href{https://github.com/EvoLandEco}{EvoLandEco}\sep
  \contacticon{linkedin}\href{https://www.linkedin.com/in/tjqin}{tjqin}\par
  \vspace{2pt}\contacticon{orcid}\href{https://orcid.org/0000-0003-2540-1778}{0000-0003-2540-1778}
}\par
\vspace{10pt}{\color{Accent}\hrule height 1.6pt}\vspace{8pt}
{\fontsize{10.1}{14}\selectfont
  \textbf{Computational biology}\sep\textbf{AI/ML}\sep\textbf{Infectious disease modelling}\par
  Network science\sep Phylodynamics\sep One Health\sep Research software
}\par
\vspace{8pt}
Computational biologist developing machine learning methods, stochastic models and research software
for biological systems. My research spans neural inference from phylogenies, dynamic livestock networks
and infectious disease modelling. I translate empirical data, scientific questions and stakeholder needs
into transparent models, reproducible HPC workflows and interactive tools for analysis and decision support.

\sectionhead{Selected research \& software}
\proof{EvoNN \& EVE: neural inference for evolutionary biology}{AI/ML \middotsep Stochastic modelling}
{Designed ensemble neural networks that learn from tree structure, branching times and summary
statistics, alongside stochastic diversification models and high-performance simulators.}
{First-author articles: \href{https://doi.org/10.1093/sysbio/syaf060}{\emph{Systematic Biology} (2026)}
\sep \href{https://doi.org/10.1016/j.jtbi.2024.111992}{\emph{Journal of Theoretical Biology} (2025)}}

\proof{NetForge: synthetic livestock contact networks}{Network science \middotsep Epidemiology}
{Developed within the IMBIT pandemic preparedness programme, combining stochastic block models
with contact history selection to generate synthetic livestock contact networks; benchmarked them
against temporal reachability and simulated infection dynamics.}
{First-author \href{https://www.biorxiv.org/content/10.64898/2026.09.15.751711v1}{\emph{bioRxiv} preprint (2026)}
\sep \href{https://qtj.me/NetForge}{Software and documentation}}

\begin{showcase}
  \smallcapslabel{Featured application}\hfill
  {\fontsize{9}{11}\selectfont\contacticon{globe}\href{https://herdlink.nl}{HerdLink.nl}\sep
  \href{https://github.com/EvoLandEco/herdlink-web}{Source code}}\par
  \vspace{4pt}
  {\fontsize{13}{16}\selectfont\bfseries HerdLink: interactive intervention modelling}\par
  \vspace{4pt}
  Developed within IMBIT and deployed as a public web platform linking Dutch livestock trade
  networks, geographic views and disease simulations. Implemented dated route and regional movement
  restrictions, seven scenario presets, and paired baseline and intervention comparisons of epidemic
  trajectories, regional burden and retained trade for scenario comparison and decision support.\par
  \vspace{4pt}
  {\fontsize{9}{11}\selectfont\color{Muted}JavaScript\middotsep React\middotsep D3.js\middotsep Vite\middotsep Automated tests and deployment}
\end{showcase}

\sectionhead{Research experience}
\datedentry{2025--present}{Postdoctoral Researcher}
{Wageningen University \& Research\middotsep Infectious Disease Epidemiology Group}
{Develop AI and computational methods for One Health and pandemic preparedness, integrating multiplex temporal networks, network-phylodynamic inference, spatio-temporal epidemic modelling, and neural/hybrid AI for outbreak detection and projection. Build reproducible HPC pipelines and interactive web tools for simulation, scenario comparison and decision support.}

\datedentry{2019--2026}{Doctoral Researcher in Theoretical \& Computational Biology}
{University of Groningen\middotsep Groningen Institute for Evolutionary Life Sciences}
{Developed simulation and inference methods for evolutionary diversification; built research
packages in R, C++ and Python with CPU/GPU computing, testing, documentation and open releases.}

\vspace{1pt}
{\fontsize{9.2}{12}\selectfont\color{Warm}\textcolor{Accent}{\faIcon{trophy}}\hspace{0.45em}\textbf{Denise Kirschner Best Student Paper Prize, 2025}
\hfill\href{https://www.sciencedirect.com/journal/journal-of-theoretical-biology/about/news/the-winner-of-the-denise-kirschner-best-student-paper-prize}{\emph{Journal of Theoretical Biology}}}

\vspace{9pt}{\color{Rule}\hrule height 0.5pt}\vspace{6pt}
{\fontsize{9.3}{12}\selectfont\textcolor{Accent}{\faIcon{language}}\hspace{0.45em}\textbf{Languages}\enspace
Chinese: native\sep English: proficient\sep Dutch: A2 certified, strong reading\sep German: A2}

\newpage
\pageheading{Research record}{Publications \& education}

\sectionhead{Peer-reviewed articles}
{\fontsize{9.4}{12.0}\selectfont
\begin{pubs}
  \item \qname, van Benthem, K. J., Valente, L.\senior, \& Etienne, R. S.\senior\ (2026).
  Parameter estimation from phylogenetic trees using neural networks and ensemble learning.
  \emph{Systematic Biology}, 75(2), 344--365. \doilink{10.1093/sysbio/syaf060}

  \item \qname, Valente, L.\senior, \& Etienne, R. S.\senior\ (2025).
  Impact of evolutionary relatedness on species diversification and tree shape.
  \emph{Journal of Theoretical Biology}, 598, 111992. \doilink{10.1016/j.jtbi.2024.111992}

  \item Sun, K., Liu, X.-S., \qjname, Jiang, F., Cai, J.-F., Shen, Y.-L., A, S.-H., \& Li, H.-L. (2022).
  Relative abundance of invasive plants more effectively explains the response of wetland communities
  to different invasion degrees than phylogenetic evenness.
  \emph{Journal of Plant Ecology}, 15(3), 625--638. \doilink{10.1093/jpe/rtab074}

  \item \qname, Zhou, J., Sun, Y., M\"uller-Sch\"arer, H., Luo, F., Dong, B., Li, H., \& Yu, F.-H. (2020).
  Phylogenetic diversity is a better predictor of wetland community resistance to
  \emph{Alternanthera philoxeroides} invasion than species richness.
  \emph{Plant Biology}, 22(4), 591--599. \doilink{10.1111/plb.13101}

  \item Wan, J.-Z., Wang, M.-Z., \qjname, Bu, X.-Q., Li, H.-L., \& Yu, F.-H. (2019).
  Spatial environmental heterogeneity may drive functional trait variation in
  \emph{Hydrocotyle vulgaris} (Araliaceae), an invasive aquatic plant.
  \emph{Aquatic Biology}, 28, 149--158. \doilink{10.3354/ab00716}

  \item \qjname*, Guan, Y.-T.*, Quan, H., Dong, B.-C., Luo, F.-L., Zhang, M.-X., Li, H.-L., \& Yu, F.-H. (2019).
  Growth traits of the exotic plant \emph{Hydrocotyle vulgaris} and the evenness of resident plant
  communities are mediated by community age, not species diversity.
  \emph{Weed Research}, 59(5), 377--386. \doilink{10.1111/wre.12373}

  \item \qjname, Guan, Y.-T., Zhang, M.-X., Li, H.-L., \& Yu, F.-H. (2018).
  Sediment type and nitrogen deposition affect the relationship between \emph{Alternanthera philoxeroides}
  and experimental wetland plant communities.
  \emph{Marine and Freshwater Research}, 69(5), 811--822. \doilink{10.1071/MF17335}

  \item Liu, L., Guan, Y.-T., \qjname, Wang, Y.-Y., Li, H.-L., \& Zhi, Y.-B. (2018).
  Effects of water regime on the growth of the submerged macrophyte \emph{Ceratophyllum demersum}
  at different densities.
  \emph{Journal of Freshwater Ecology}, 33(1), 45--56. \doilink{10.1080/02705060.2017.1422043}
\end{pubs}
}

\sectionhead{Preprints}
{\fontsize{9.4}{12.0}\selectfont
\begin{pubs}
  \item \qname, Atamer Balkan, B., Schmid, B. V., \& ten Bosch, Q. A. (2026).
  \textbf{From Movement to Spread: Generating Livestock Contact Networks that Preserve Infection Dynamics.}
  \emph{bioRxiv}. \doilink{10.64898/2026.09.15.751711}

  \item \qname, van Benthem, K., Valente, L.\senior, \& Etienne, R. S.\senior\ (2026).
  Identifying evolutionary relatedness effects on diversification from phylogenies using neural networks.
  \emph{bioRxiv}. \doilink{10.64898/2026.01.12.699140}
\end{pubs}
}
\vspace{3pt}{\fontsize{8.3}{10}\selectfont\color{Muted}* Equal contribution. \senior\ Joint senior authors.}\par

\sectionhead{Conference presentation}
{\fontsize{9.4}{12.0}\selectfont
\qname, Atamer Balkan, B., Schmid, B., \& ten Bosch, Q. (2026).
\href{https://modahconf2026.sciencesconf.org/program?lang=en}{From Livestock Movements to Infection Dynamics: Mechanism-Informed Learning for Synthetic Contact Networks}.
\textcolor{Accent}{\faIcon{microphone}}\hspace{0.35em}\textbf{Oral presentation},
\emph{ModAH Conference 2026}, Nantes, France, 26 August 2026.
}\par

\sectionhead{Education}
{\fontsize{9.4}{12.0}\selectfont
\datedentry{2019--2026}{PhD in Theoretical Biology}{University of Groningen, Netherlands}
{Thesis: \href{https://research.rug.nl/files/1515277445/Complete_thesis.pdf}{\emph{Diversification Models and Neural Inference}}.
Supervisors: Rampal Etienne, Luis Valente and Koen van Benthem.}
\datedentry{2016--2019}{MSc in Wetland Ecology}{Beijing Forestry University, China}
{Supervisors: Hongli Li and Feihai Yu. Research on wetland invasion resistance and community
assembly, with field surveys across 18 Chinese provinces, controlled experiments, image analysis
and phylogenetic workflows.}
\datedentry{2012--2016}{BSc in Marine Biology}{Nanjing Normal University, China}{}
}

\newpage
\pageheading{Research practice}{Programmes, software \& expertise}

\compactsectionhead{Research programmes \& funding}
\programme{NextdAI: Integrated AI-Driven Early Detection System}{In development}{2026}
{NWO KIC: Early Detection of (Emerging) Infectious Disease Outbreaks}
{Co-develop AI architecture and data infrastructure for heterogeneous surveillance, anomaly detection,
neurosymbolic characterisation, forecasting and decision support.}

\programme{EUPAHW: hepatitis E across pigs, wild boar and humans}{Ongoing}{2026}
{European Partnership on Animal Health and Welfare\middotsep SOA26}
{Develop phylodynamic and network analyses linking livestock movements, wildlife and landscape
data to investigate transmission pathways and inform One Health surveillance.}

\programme{SSS-mod: Preparedness Package for Future Respiratory Outbreaks}{Ongoing}{2026}
{ZonMw Implementatie-impuls COVID-19 programma}
{Develop data preparation, recalibration, model output and interface workflows supporting reuse of a
setting-specific respiratory transmission model.}

\programme{IMBIT: Intervention Modeling \& Bayesian Inference of Transmission Routes for Early Pandemic Response}{Completed}{2025}
{ZonMw Kennisprogramma Pandemische Paraatheid}
{Developed NetForge and HerdLink.nl for synthetic livestock contact networks and interactive
intervention experiments. Contributed to a One Health programme coupling multi-route simulation with
Bayesian inference of transmission routes across human, livestock, wildlife and environmental systems.}

\compactsectionhead{Research software}
{\fontsize{9.35}{11.5}\selectfont
\begin{tabularx}{\linewidth}{@{}>{\RaggedRight\arraybackslash}p{3.55cm}@{\hspace{0.35cm}}Y@{}}
\skillrow{\href{https://github.com/EvoLandEco/EvoNN}{EvoNN} \& \href{https://github.com/EvoLandEco/evesim}{evesim}}
{Developer. Neural inference from phylogenies and fast simulation of evolutionary relatedness dependent birth--death processes.}
\skillrow{\href{https://github.com/EvoLandEco/NetSpectra}{NetSpectra}, \href{https://github.com/EvoLandEco/EvoLab}{EvoLab}, \href{https://github.com/EvoLandEco/miniape}{miniape}}
{Developer. Interactive network dynamics, lineage simulation, tree reconstruction and browser-based phylogenetic data handling.}
\skillrow{\href{https://github.com/thijsjanzen/treestats}{treestats}, \href{https://github.com/rsetienne/DDD}{DDD}, \href{https://github.com/rsetienne/DAISIE}{DAISIE}}
{Collaborator. Phylogenetic statistics, data structures, simulation and inference for diversification and island biogeography.}
\end{tabularx}
}

\sectionhead{Technical expertise}
{\fontsize{9.7}{12.5}\selectfont
\begin{tabularx}{\linewidth}{@{}>{\RaggedRight\arraybackslash}p{3.7cm}@{\hspace{0.4cm}}Y@{}}
\skillrow{AI \& machine learning}{Graph, recurrent and dense neural networks; Transformers; ensemble and generative modelling;
PyTorch, PyTorch Geometric, model benchmarking and evaluation.}
\skillrow{Modelling \& inference}{Stochastic processes, simulation-based inference, likelihood methods,
dynamic and temporal networks, transmission simulation, phylodynamics and phylogenetic analysis.}
\skillrow{Programming \& data}{Python, R, C/C++, JavaScript, SQL; package development,
spatiotemporal data integration, GIS and reproducible analysis pipelines.}
\skillrow{Research computing}{Linux, Bash, CPU/GPU HPC, Docker, Git, GitHub workflows,
CI/CD, automated testing, reproducible environments and AI-assisted software development.}
\skillrow{Web \& visualisation}{React, D3.js, Three.js, Vite; interactive scientific interfaces,
geospatial visualisation, \LaTeX, TikZ and PGFPlots.}
\end{tabularx}
}

\compactsectionhead{Seminars, workshops, mentoring \& service}
{\fontsize{9.25}{11.2}\selectfont
\begin{cvitems}
  \item Seminars and workshops at Groningen, WUR and Wageningen Bioveterinary Research on island biogeography,
  stochastic modelling, phylodynamics, neural networks and network statistics.
  \item Mentor BSc/MSc students and junior PhD researchers; pair programming on test-driven development
  and scientific software design.
  \item Reviewer for \emph{Systematic Biology}, \emph{PLOS Computational Biology}, \emph{Epidemics},
  \emph{BMC Public Health} and \emph{BMC Plant Biology}; open-source contributor.
\end{cvitems}
\vspace{1.5pt}
{\fontsize{9.25}{11.2}\selectfont
\textbf{Academic leadership:} Vice President, Graduate Student Council (2017--2018); student representative,
Wetland Ecology Group (2016--2019). \textbf{University awards:} First-Class Academic Performance Award
(2017--2019); Graduate Academic Innovation Award (2018), Beijing Forestry University.}
}\par

\end{document}
